\setcounter{section}{24}

Misalkan

\mavergleichskette
{\vergleichskette
{ Y
}
{ \subseteq }{ \R^n
}
{  }{
}
{    }{
}
{   }{
}
}
{}{}{ }
adalah suatu hipermuka terdiferensialkan. Bagi setiap titik

\mavergleichskette
{\vergleichskette
{ P
}
{ \in }{ Y
}
{  }{
}
{    }{
}
{   }{
}
}
{}{}{,}
\definitionsverweis {ruang singgung}{}{}
merupakan subruang vektor

\mavergleichskette
{\vergleichskette
{ T_PY
}
{ \subseteq }{ \R^n
}
{  }{
}
{    }{
}
{   }{
}
}
{}{}{}
di ruang ambien. Dengan menggunakan hasil kali dalam standar pada $\R^n$, kita memperoleh
\definitionsverweis {proyeksi ortogonal}{}{}
\maabbdisp {\pi_P} {\R^n } { T_PY
} {,}
yang memungkinkan data di ruang ambien ditafsirkan sebagai data tangensial pada hipermuka. Dari konstruksi ini muncul berbagai konsep penting pada $Y$, seperti
\definitionsverweis {percepatan tangensial}{}{,}
\definitionsverweis {kurva geodesik}{}{,}
\definitionsverweis {kelengkungan}{}{,}
\definitionsverweis {turunan kovarian suatu medan vektor}{}{,}
\definitionsverweis {medan vektor paralel}{}{}
serta
\definitionsverweis {transpor paralel}{}{.}
Di sisi lain, hasil kali dalam standar tersebut juga menginduksi
\definitionsverweis {struktur Riemann}{}{}
pada $Y$. Pertanyaannya kemudian ialah apakah konsep-konsep tadi dapat dibangun secara intrinsik hanya dari $Y$ dan struktur Riemann-nya
\zusatzklammer {misalnya untuk
\definitionsverweis {permukaan hiperbolik}{}{}} {} {.}
Untuk menjawabnya, diperlukan perangkat baru, yaitu koneksi pada bundel vektor, terutama pada bundel tangen suatu
\definitionsverweis {manifold diferensiabel}{}{.}
Pada manifold Riemann, struktur ini menghasilkan secara kanonik \stichwort {koneksi Levi-Civita} {}
\zusatzklammer {lihat
Teorema 26.12} {} {,}
yang memungkinkan seluruh konsep di atas dirumuskan secara intrinsik.

  \zwischenueberschrift{Koneksi}

Misalkan $X$ adalah
\definitionsverweis {manifold diferensiabel}{}{}
riil, dan
\maabbdisp {p} { E } { X
} {}
adalah bundel vektor di atas $X$ yang ruang totalnya juga diperlakukan sebagai manifold diferensiabel. Adapun
\definitionsverweis {pemetaan singgung}{}{}
\maabb {T(p)} {TE} {TX
} {,}
berada dalam diagram komutatif

\mathdisp {\begin{matrix} TE & \stackrel{ T(p) }{\longrightarrow} & TX &  \\ \downarrow & & \downarrow  & \\  E  & \stackrel{ p }{\longrightarrow} &  X  & \! \! \! \!\!\!\!\!   \\ \end{matrix}} {  }

dan, bagi setiap titik

\mavergleichskette
{\vergleichskette
{ e
}
{ \in }{ E
}
{  }{
}
{    }{
}
{   }{
}
}
{}{}{}
dengan titik dasar

\mavergleichskette
{\vergleichskette
{ x
}
{ = }{ p(e)
}
{  }{
}
{    }{
}
{   }{
}
}
{}{}{,}
\definitionsverweis {pemetaan singgung}{}{}
\definitionsverweis {linear}{}{}
\maabbdisp {T_e(p)} {T_e E} { T_{x}X
} {}
bersifat surjektif. Salah satu alasannya ialah bahwa secara lokal, melalui setiap titik

\mavergleichskette
{\vergleichskette
{ e
}
{ \in }{ E
}
{  }{
}
{    }{
}
{   }{
}
}
{}{}{,}
terdapat suatu penampang bundel vektor $E$. Agar domain dan kodomain memiliki ruang dasar yang sama, kita memandang $T(p)$ sebagai
\definitionsverweis {pemetaan bundel}{}{}
\maabbdisp {Tp} {TE} { p^* TX
} {}
antara bundel vektor di atas $E$. Dalam notasi ini,

\mavergleichskettedisp
{\vergleichskette
{ p^*TX
}
{ =} { E \oplus TX
}
{ =} { E \times_X TX
}
{ } {
}
{ } {
}
}
{}{}{}
merupakan
\definitionsverweis {jumlah langsung}{}{}
bundel vektor di atas $X$; lihat
Soal 12.24.
Di sini, $p^*TX$ dipandang sebagai bundel vektor di atas $E$. Serat $p^*TX$ di atas titik

\mavergleichskette
{\vergleichskette
{ e
}
{ \in }{ E
}
{  }{
}
{    }{
}
{   }{
}
}
{}{}{}
dengan titik dasar

\mavergleichskette
{\vergleichskette
{ x
}
{ = }{ p(e)
}
{  }{
}
{    }{
}
{   }{
}
}
{}{}{}
hanyalah $T_{x}X$, sedangkan pemetaan pada serat tersebut tetaplah
\maabbdisp {T_e (p)} { T_e E } { T_{x} X
} {,}
sehingga di dalam $p^*TX$ terdapat satu salinan $T_{x} X$ untuk setiap

\mavergleichskette
{\vergleichskette
{ e
}
{ \in }{ E
}
{  }{
}
{    }{
}
{   }{
}
}
{}{}{.}

\inputbemerkung
{}
{

Misalkan
\maabb {p} { E } { X
} {}
adalah
\definitionsverweis {bundel vektor}{}{}
diferensiabel di atas
\definitionsverweis {manifold diferensiabel}{}{}
$X$. Ambil himpunan terbuka

\mavergleichskette
{\vergleichskette
{ U
}
{ \subseteq }{ X
}
{  }{
}
{    }{
}
{   }{
}
}
{}{}{,}
tempat $E$ dapat ditrivialkan, sehingga

\mavergleichskette
{\vergleichskette
{ E  \vert_U
}
{ \cong }{ U \times W
}
{  }{
}
{    }{
}
{   }{
}
}
{}{}{}
dengan ruang vektor atas $\R$ yang kita namakan $W$. Dengan demikian,

\mavergleichskettedisp
{\vergleichskette
{ T  \left(  E \vert_U  \right)
}
{ =} { T(U \times  W )
}
{ =} { T U \times T W
}
{ =} { TU \times W \times W
}
{ } {
}
}
{}{}{,}
Sementara itu,
\definitionsverweis {pemetaan singgung}{}{}
yang berkaitan dengan
\maabb {p_U} {E \vert_U } { U
} {}
merupakan proyeksi ke $TU$. Untuk menuliskannya sebagai homomorfisme bundel di atas $U\times W$, kodomainnya diidentifikasi sebagai

\mavergleichskettedisp
{\vergleichskette
{ p^*TU
}
{ =} { TU \times W
}
{ =} { TU \times_U (U \times W)
}
{ } {
}
{ } {
}
}
{}{}{}
dan pemetaannya ialah
\maabbdisp {} {T  \left(  E \vert_U  \right)  = TU \times W \times W  } { TU \times W
} {}
yang mengambil proyeksi ke $TU$ sekaligus proyeksi kedua ke $W$.

}

\inputdefinition
{}
{

Misalkan
\maabb {p} { E } { X
} {}
adalah
\definitionsverweis {bundel vektor diferensiabel}{}{}
di atas
\definitionsverweis {manifold diferensiabel}{}{}
$X$.
\definitionsverweis {Bundel kernel}{}{}
dari
\definitionsverweis {homomorfisme bundel}{}{}
\definitionsverweis {surjektif}{}{}
\maabbdisp {T(p)} {TE} {p^*TX
} {}
di atas $E$ disebut \definitionswort {bundel vertikal}{.} Notasinya adalah $V$.

}

Menurut Soal 12.19, bundel vertikal adalah bundel vektor di atas $E$; lebih khusus lagi, ia merupakan
\definitionsverweis {subbundel vektor}{}{}
dari $TE$.



\inputfaktbeweis
{Mannigfaltigkeit/Vektorbündel/R/Vertikalbündel/Eigenschaften/Fakt}
{Lema}
{}
{

\faktsituation {Misalkan
\maabb {p} { E } { X
} {}
adalah
\definitionsverweis {bundel vektor diferensiabel}{}{}
di atas
$C^2$-\definitionsverweis {manifold diferensiabel}{}{}
$X$.}
\faktuebergang {Pernyataan berikut berlaku.}
\faktfolgerung {\aufzaehlungdrei{Terdapat
\definitionsverweis {barisan eksak pendek}{}{}

\mathdisp {0 \, \stackrel{  }{\longrightarrow} \,  V   \, \stackrel{  }{ \longrightarrow}   \,  TE  \, \stackrel{  }{\longrightarrow} \, p^* TX \,\stackrel{  } {\longrightarrow}  \, 0} {  }

yang terdiri dari
\definitionsverweis {homomorfisme}{}{}
bundel vektor di atas $E$.
}{Untuk
\definitionsverweis {bundel vertikal}{}{}
terdapat
\definitionsverweis {isomorfisme bundel}{}{}

\mavergleichskette
{\vergleichskette
{ V \cong p^*E
}
{ = }{ E \oplus E
}
{  }{
}
{    }{
}
{   }{
}
}
{}{}{.}
}{Bagi setiap titik

\mavergleichskette
{\vergleichskette
{ e
}
{ \in }{ E
}
{  }{
}
{    }{
}
{   }{
}
}
{}{}{,}
diperoleh barisan eksak pendek ruang vektor berikut:

\mathdisp {0 \, \stackrel{  }{\longrightarrow} \, E_{p(e)}   \, \stackrel{  }{ \longrightarrow}   \,  T_eE  \, \stackrel{  }{\longrightarrow} \, T_{p(e)}X \,\stackrel{  } {\longrightarrow}  \, 0} {  }

}}
\faktzusatz {}
\faktzusatz {}

}
{

\aufzaehlungzwei {Pernyataan ini langsung mengikuti definisi.
} {Terdapat homomorfisme injektif bundel vektor
\maabbdisp {} {p^*E = E \oplus E} { TE
} {}
di atas $E$ yang memetakan
\mathl{(P,u)}{,}

\mavergleichskette
{\vergleichskette
{ u
}
{ \in }{ E_P
}
{  }{
}
{    }{
}
{   }{
}
}
{}{}{,}
ke vektor singgung yang diwakili oleh kurva diferensiabel
\zusatzklammer {vertikal} {} {}
\maabbeledisp {} {\R} {E
} { t } {(P,tu)
} {,}
ini. Selanjutnya, oleh pemetaan
\maabbdisp {} {TE} { p^*TX
} {}
vektor singgung tersebut dikirim ke $0$, sebab kurva wakilnya dikirim ke $P$ oleh
\maabb {p} { E } { X
} {}
.
Oleh karena itu,

\mavergleichskette
{\vergleichskette
{ E
}
{ \subseteq }{ V
}
{  }{
}
{    }{
}
{   }{
}
}
{}{}{,}
dan perbandingan rank menunjukkan bahwa inklusi ini sebenarnya suatu isomorfisme.
}

}

Kasus yang sering dijumpai adalah ketika $E$ merupakan bundel tangen dari $X$. Dalam kasus ini diperoleh barisan eksak pendek

\mathdisp {0 \, \stackrel{  }{\longrightarrow} \, p^* TX  \, \stackrel{  }{ \longrightarrow}   \, TTX  \, \stackrel{  }{\longrightarrow} \, p^* TX \,\stackrel{  } {\longrightarrow}  \, 0} {  }

Oleh sebab itu, kedua salinan bundel tangen di ruas kiri dan kanan harus selalu dibedakan dengan cermat.

\inputbemerkung
{}
{

Misalkan
\maabbdisp {f} {\R^n } { \R
} {}
adalah fungsi yang dua kali
\definitionsverweis {terdiferensialkan secara kontinu}{}{,}
dan

\mavergleichskette
{\vergleichskette
{ Y
}
{ = }{ f^{-1} (0)
}
{  }{
}
{    }{
}
{   }{
}
}
{}{}{}
adalah
\definitionsverweis {hipermuka terdiferensialkan}{}{}
terkait yang
\definitionsverweis {reguler}{}{}
pada setiap titik.
\definitionsverweis {Bundel tangen}{}{}
$TY$ dapat pula dinyatakan sebagai serat bersama dari fungsi-fungsi, yakni

\mavergleichskettealignhandlinks
{\vergleichskettealignhandlinks
{ TY
}
{ =} { \left\{  (P,u)   \mid   f(P) = 0 , \,  \left(  D f   \right)_{ P } \left(  u  \right) = 0       \right\}
}
{ =} { \left\{  (P,u_1 , \ldots , u_n)   \mid   f(P) = 0 , \,  \sum_{i = 1}^n u_i { \frac{   \partial   f   }{  \partial   x_i   } } ( P) = 0       \right\}
}
{ \subseteq} { \R^n \times \R^n
}
{ } {
}
}
{}
{}{,}
sebagaimana dalam
Soal 10.15.
Dalam rumusan ini, $f$ dipandang sebagai fungsi pada $\R^{2n}$ yang hanya bergantung pada $n$ variabel pertama, dan kita definisikan

\mavergleichskettedisp
{\vergleichskette
{ g(P,u)
}
{ =} { \sum_{i = 1}^n u_i { \frac{   \partial   f   }{  \partial   x_i   } } ( P)
}
{ } {
}
{ } {
}
{ } {
}
}
{}{}{.}
Bundel tangen kedua $TTY$, yaitu bundel tangen dari $TY$, selanjutnya dapat dideskripsikan dengan
\definitionsverweis {diferensial total}{}{}
dari pemetaan
\maabbdisp {\varphi = (f,g)} {\R^{2n}} { \R^2
} {}
sebagai berikut:

\mavergleichskettealignhandlinks
{\vergleichskettealignhandlinks
{ TTY
}
{ =} { \left\{  (P,u, v,w)  \mid   (P,u) \in TY , \,  \left(  D\varphi  \right)_{(P,u)} \left(  (v,w)  \right) = 0       \right\}
}
{ \subseteq} { \R^{4n}
}
{ } {
}
{ } {}
}
{}
{}{.}
\definitionsverweis {Matriks Jacobi}{}{,}
yang merepresentasikan diferensial total $D\varphi$, ialah

\mathdisp {\begin{pmatrix}   { \frac{   \partial   f   }{  \partial   x_1   } } ( P)   &   \ldots &  { \frac{   \partial   f   }{  \partial   x_n   } } ( P)  &   0 &  \ldots &  0  \\   \sum_{i = 1}^n u_i   { \frac{   \partial    }{  \partial   x_1   } } { \frac{   \partial   f   }{  \partial   x_i   } } ( P)  &  \ldots  &   \sum_{i = 1}^n u_i   { \frac{   \partial    }{  \partial   x_n   } } { \frac{   \partial   f   }{  \partial   x_i   } } ( P)   &  { \frac{   \partial   f   }{  \partial   x_1   } } ( P)  &  \ldots &  { \frac{   \partial   f   }{  \partial   x_n   } } ( P)    \end{pmatrix}} { . }

Di samping dua syarat yang mendefinisikan $TY$ dan hanya menyangkut
\mathkor {} {P} {dan} {u} {,}
diperoleh dua syarat tambahan yang melibatkan seluruh variabel, yaitu

\mavergleichskettedisp
{\vergleichskette
{ \sum_{j = 1}^n v_j { \frac{   \partial   f   }{  \partial   x_j   } } ( P)
}
{ =} { 0
}
{ } {
}
{ } {
}
{ } {
}
}
{}{}{}
dan

\mavergleichskettedisp
{\vergleichskette
{ \sum_{j = 1}^n v_j  \left(  \sum_{i = 1}^n u_i { \frac{   \partial    }{  \partial   x_i   } }  { \frac{   \partial   f   }{  \partial   x_j   } } ( P)  \right)      + \sum_{k = 1}^n w_k { \frac{   \partial   f   }{  \partial   x_k   } } ( P)
}
{ =} { 0
}
{ } {
}
{ } {
}
{ } {
}
}
{}{}{.}

}

\inputbemerkung
{}
{

Misalkan
\maabbdisp {f} {\R^n } { \R
} {}
adalah fungsi yang dua kali
\definitionsverweis {terdiferensialkan secara kontinu}{}{,}
dan

\mavergleichskette
{\vergleichskette
{ Y
}
{ = }{ f^{-1} (0)
}
{  }{
}
{    }{
}
{   }{
}
}
{}{}{}
adalah
\definitionsverweis {hipermuka terdiferensialkan}{}{}
terkait yang
\definitionsverweis {reguler}{}{}
pada setiap titik. Kita gunakan deskripsi bundel tangen kedua $TTY$ dari
Catatan 24.4
beserta notasi di sana. Bundel tangen tarik balik $\pi^*TY$ terhadap
\maabbdisp {\pi} {TY} {Y
} {,}
yang juga merupakan
\definitionsverweis {bundel vertikal}{}{,}
adalah

\mavergleichskettealign
{\vergleichskettealign
{ V
}
{ \cong} { \pi^*TY
}
{ =} { TY \times_Y TY
}
{ =} { \left\{ (P,u,z)  \mid  f(P) = 0 , \,  \sum_{i = 1}^n u_i { \frac{   \partial   f   }{  \partial   x_i   } } ( P) = 0   , \,  \sum_{k = 1}^n z_k { \frac{   \partial   f   }{  \partial   x_k   } } ( P) = 0    \right\}
}
{ } {
}
}
{}
{}{.}
Sekarang kita tentukan bentuk pemetaan singgung
\maabb {T(\pi)} {TTY } { \pi^* TY
} {}
serta pembenaman bundel vertikal ke dalam $TTY$. Berdasarkan
Soal 24.2,
diperoleh
\maabbeledisp {T(\pi)} {TTY} { \pi^* TY
} {(P,u,v,w)} { (P,u,v)
} {.}
Kernel pemetaan ini adalah

\mavergleichskettedisp
{\vergleichskette
{ V
}
{ =} { \left\{  (P,u,0,w)   \mid   (P,u,0,w) \in TTY          \right\}
}
{ \subset} { TTY
}
{ } {
}
{ } {
}
}
{}{}{,}
dan langsung terlihat bahwa
\mathl{(P,u,w)}{} memenuhi syarat bagi
\mathl{\pi^*TY}{}.
Oleh sebab itu, kedua homomorfisme bundel dalam barisan eksak pendek

\mathdisp {0 \, \stackrel{  }{\longrightarrow} \,  \pi^*TY  \, \stackrel{  }{ \longrightarrow}   \, TTY  \, \stackrel{  }{\longrightarrow} \,  \pi^*TY \,\stackrel{  } {\longrightarrow}  \, 0} {  }

dapat ditulis dalam koordinat sebagai

\mathl{(P,u,w) \mapsto  (P,u,0,w)}{} dan
\mathl{(P,u,v,w) \mapsto (P,u,v)}{}.

}

Jika $E$ merupakan bundel trivial, yaitu

\mavergleichskette
{\vergleichskette
{ E
}
{ = }{ X \times \R^r
}
{  }{
}
{    }{
}
{   }{
}
}
{}{}{,}
maka terdapat penguraian langsung
\zusatzklammer {di atas $E$} {} {}

\mavergleichskette
{\vergleichskette
{ TE
}
{ = }{ V \oplus p^* TX
}
{  }{
}
{    }{
}
{   }{
}
}
{}{}{.}
Di setiap titik

\mavergleichskette
{\vergleichskette
{ e
}
{ \in }{ E
}
{  }{
}
{    }{
}
{   }{
}
}
{}{}{,}
vektor singgung

\mavergleichskette
{\vergleichskette
{ t
}
{ \in }{ V
}
{  }{
}
{    }{
}
{   }{
}
}
{}{}{}
mewakili arah vertikal, sedangkan

\mavergleichskette
{\vergleichskette
{ t
}
{ \in }{ p^* T_xX
}
{  }{
}
{    }{
}
{   }{
}
}
{}{}{}
mewakili \anfuehrung{arah horizontal}{.}
Pada bundel vektor sembarang, suatu trivialisasi lokal menguraikan ruang singgung
\mathl{T_eE}{} sebagai jumlah langsung ruang vertikal dan ruang horizontal. Namun, subruang horizontal yang dihasilkan sangat bergantung pada trivialisasi yang dipilih. Tidak seperti subruang vertikal
\zusatzklammer {yang bersama-sama membentuk bundel vertikal} {} {,}
subruang horizontal ini belum tentu menyatu menjadi suatu subbundel. Keberadaan subbundel horizontal semacam inilah yang dirumuskan oleh konsep koneksi.

\inputdefinition
{}
{

Misalkan $X$ adalah
\definitionsverweis {manifold diferensiabel}{}{}
dan $E$ adalah
\definitionsverweis {bundel vektor diferensiabel}{}{}
di atas $X$. Suatu
\definitionswort {koneksi}{}
pada $E$ didefinisikan sebagai penguraian jumlah langsung
\definitionsverweis {bundel tangen}{}{}

\mavergleichskette
{\vergleichskette
{ TE
}
{ = }{ V \oplus H
}
{  }{
}
{    }{
}
{   }{
}
}
{}{}{}
menjadi dua
\definitionsverweis {subbundel vektor}{}{}
\mathkor {} {V} {dan} {H} {,}
dengan ketentuan bahwa

\mavergleichskette
{\vergleichskette
{ V
}
{ \subseteq }{ TE
}
{  }{
}
{    }{
}
{   }{
}
}
{}{}{}
merupakan
\definitionsverweis {bundel vertikal}{}{.}
Subbundel

\mavergleichskette
{\vergleichskette
{ H
}
{ \subseteq }{ TE
}
{  }{
}
{    }{
}
{   }{
}
}
{}{}{}
disebut
\definitionswort {bundel horizontal}{}
koneksi tersebut.

}

Koneksi biasanya dilambangkan dengan $\nabla$, khususnya ketika proses diferensiasi yang ditentukannya hendak ditekankan. Jadi, menurut definisi, koneksi hanyalah subbundel dari $TE$ yang komplementer terhadap bundel vertikal. Ada beberapa rumusan lain yang ekuivalen. Memberikan suatu koneksi sama artinya dengan memberikan homomorfisme bundel
\zusatzklammer {disebut \stichwort {proyeksi vertikal} {} koneksi} {} {}
\maabbdisp {\pi_{\text{vert} }} {TE} {V
} {}
dengan

\mavergleichskettedisp
{\vergleichskette
{ \pi_{\text{vert} }  \circ \iota
}
{ =} {
\operatorname{Id}_{  V  }
}
{ } {
}
{ } {
}
{ } {
}
}
{}{}{,}
di mana $\iota$ adalah pembenaman $V$ ke dalam $TE$. Dalam rumusan ini, bundel horizontal ialah kernel $\pi_{\text{vert}}$. Sebaliknya, penguraian jumlah langsung menyediakan proyeksi ke masing-masing suku. Penguraian seperti itu juga disebut \stichwort {pemisahan} {} dari barisan eksak pendek. Setiap koneksi dapat dibatasi ke himpunan terbuka

\mavergleichskette
{\vergleichskette
{ U
}
{ \subseteq }{ X
}
{  }{
}
{    }{
}
{   }{
}
}
{}{}{.}

\inputbeispiel{}
{

Pertimbangkan bundel trivial
\maabbdisp {p} {X \times W = W_X } { X
} {,}
dengan $W$ ruang vektor riil
\definitionsverweis {berdimensi}{}{}
$r$ dan $X$ sebuah
\definitionsverweis {manifold diferensiabel}{}{,}
kita memperoleh barisan eksak pendek

\mathdisp {0 \, \stackrel{  }{\longrightarrow} \, p^*(X \times W)  =  X \times W \times W   \, \stackrel{  }{ \longrightarrow}   \, p^* T(X \times W )  = TX \times W \times W   \, \stackrel{  }{\longrightarrow} \, p^* TX  =  TX \times W   \,\stackrel{  } {\longrightarrow}  \, 0} {  }

bundel vektor di atas
\mathl{X \times W}{.}
Semua produk di sini dipahami sebagai produk Kartesius biasa, setelah dilakukan identifikasi

\mavergleichskette
{\vergleichskette
{ (X \times W) \times_X (X \times W)
}
{ = }{ X \times W \times W
}
{  }{
}
{    }{
}
{   }{
}
}
{}{}{}
di sebelah kiri serta identifikasi

\mavergleichskette
{\vergleichskette
{ TX \times_X (X \times W)
}
{ = }{ TX \times W
}
{  }{
}
{    }{
}
{   }{
}
}
{}{}{}
di sebelah kanan. Pada titik

\mavergleichskette
{\vergleichskette
{ (x,w)
}
{ \in }{ X \times W
}
{  }{
}
{    }{
}
{   }{
}
}
{}{}{,}
barisan eksak pada serat berbentuk

\mathdisp {0 \, \stackrel{  }{\longrightarrow} \,  W   \, \stackrel{  }{ \longrightarrow}   \,  T_xX \times W   \, \stackrel{  }{\longrightarrow} \,  T_xX \,\stackrel{  } {\longrightarrow}  \, 0} { . }

Pemisahan baku dari barisan eksak pendek bundel vektor ini disebut \stichwort {koneksi trivial} {} pada bundel trivial tersebut.

}

\inputbemerkung
{}
{

Misalkan
\maabbdisp {f} {\R^n } { \R
} {}
adalah fungsi yang dua kali
\definitionsverweis {terdiferensialkan secara kontinu}{}{,}
dan

\mavergleichskette
{\vergleichskette
{ Y
}
{ = }{ f^{-1} (0)
}
{  }{
}
{    }{
}
{   }{
}
}
{}{}{}
adalah
\definitionsverweis {hipermuka terdiferensialkan}{}{}
terkait yang
\definitionsverweis {reguler}{}{}
pada setiap titik. Kita kembali memakai deskripsi bundel tangen kedua $TTY$ dari
Catatan 24.4
dan homomorfisme bundel dari
Catatan 24.5.
Melalui proyeksi ortogonal
\maabb {} {\R^n } { T_PY
} {}
diperoleh proyeksi vertikal
\maabbeledisp {} {TTY} { V = p^*TY
} {(P,u,v,w)} { (P,u, \pi_P (w) )
} {,}
pada
\definitionsverweis {bundel vertikal}{}{.}
Memang, $\pi_P(w)$ berada dalam ruang singgung $Y$. Apabila kita mulai dari vektor

\mavergleichskette
{\vergleichskette
{ (P,u,w)
}
{ \in }{ V
}
{  }{
}
{    }{
}
{   }{
}
}
{}{}{,}
komponen terakhir $w$ pada
\mathl{(P,u,0,w)}{} sudah otomatis berada dalam ruang singgung. Karena itu, proyeksi ortogonal mengirim unsur ini kembali ke $(P,u,w)$. Jadi, konstruksi tersebut mendefinisikan
\definitionsverweis {koneksi}{}{}
pada $TY$.

}

  \zwischenueberschrift{Turunan vertikal}

\inputdefinition
{}
{

Misalkan $X$ adalah
\definitionsverweis {manifold diferensiabel}{}{}
dan $E$ adalah
\definitionsverweis {bundel vektor diferensiabel}{}{}
di atas $X$ yang dilengkapi dengan suatu
\definitionsverweis {koneksi}{}{.}
\definitionswort {turunan vertikal}{}
$\nabla$ yang berkaitan dengan koneksi tersebut adalah pemetaan
\maabbeledisp {\nabla} {C^1(E)} { C^0(E \otimes T^* X )
} { s } { \nabla(s) = \pi_{\text{vert} }  \circ T(s)
} {.}

}
Dengan kata lain, setiap penampang diferensiabel $s$ dari $E$
\zusatzklammer {di atas $X$ ataupun di atas himpunan terbuka sembarang

\mavergleichskettek
{\vergleichskettek
{ U
}
{ \subseteq }{ X
}
{  }{
}
{    }{
}
{   }{
}
}
{}{}{}} {} {}
dipetakan ke komposisi

\mathdisp {TX \stackrel{T(s)}{ \longrightarrow } TE \stackrel{ \pi_{\text{vert} } }{\longrightarrow} V \cong p^*E \longrightarrow E} {  }

dengan identifikasi

\mavergleichskette
{\vergleichskette
{ \operatorname{Hom}(TX,E)
}
{ \cong }{ E \otimes T^*X
}
{  }{
}
{    }{
}
{   }{
}
}
{}{}{.}

Jika diberikan pula
\definitionsverweis {medan vektor}{}{}
$F$ pada $X$, yaitu penampang
\maabb {F} { X } { TX
} {}
pada
\definitionsverweis {bundel tangen}{}{}
$TX$, maka diperoleh pemetaan
\maabbeledisp {\nabla_F} {C^1(E)} {C^0(E)
} {s } {\nabla_F(s) = \nabla (s) \circ F
} {,}
yang disebut \stichwort {turunan vertikal} {} dalam arah $F$. Ketergantungan terhadap $F$ hanya bersifat titik demi titik. Sebaliknya, ketergantungan terhadap $s$ bersifat infinitesimal karena melibatkan pemetaan singgung
\mathl{T(s)}{}.



\inputfaktbeweis
{Mannigfaltigkeit/Vektorbündel/R/Zusammenhang/Ableitung bezüglich Vektorfeld/Eigenschaften/Fakt}
{Lema}
{}
{

\faktsituation {Misalkan $X$ adalah
\definitionsverweis {manifold diferensiabel}{}{}
dan $E$ adalah
\definitionsverweis {bundel vektor diferensiabel}{}{}
di atas $X$ yang dilengkapi dengan suatu
\definitionsverweis {koneksi}{}{.}
Tuliskan
\maabbdisp {\nabla_F} {C^1(E)} {C^0(E)
} {}
untuk operator turunan yang bersesuaian dengan
\definitionsverweis {medan vektor}{}{}
$F$.}
\faktuebergang {Sifat-sifat berikut berlaku.}
\faktfolgerung {\aufzaehlungzwei {
\mathl{\left(  \nabla_F(s)  \right)  (P)}{} hanya bergantung pada $F(P)$
\zusatzklammer {serta pada $s$} {} {.}
} {Berlaku

\mavergleichskettedisp
{\vergleichskette
{ \nabla_{g_1F_1+g_2F_2 }
}
{ =} { g_1 \nabla_{F_1 } +g_2 \nabla_{F_2 }
}
{ } {
}
{ } {
}
{ } {
}
}
{}{}{}
untuk medan vektor $F_1,F_2$ dan fungsi
\maabb {g_1,g_2} { X} { \R
} {.}
}}
\faktzusatz {}
\faktzusatz {}

}
{

\aufzaehlungzwei {Tetapkan titik

\mavergleichskette
{\vergleichskette
{ P
}
{ \in }{ X
}
{  }{
}
{    }{
}
{   }{
}
}
{}{}{}
dan penampang

\mavergleichskette
{\vergleichskette
{ s
}
{ \in }{ C^1(U,E)
}
{  }{
}
{    }{
}
{   }{
}
}
{}{}{}
pada himpunan terbuka

\mavergleichskette
{\vergleichskette
{ P
}
{ \in }{ U
}
{ \subseteq }{ X
}
{    }{
}
{   }{
}
}
{}{}{,}
maka

\mavergleichskettedisp
{\vergleichskette
{ \nabla_F (s)(P)
}
{ =} { ( \pi_{\text{vert} } \circ  T_P(s) ) (F(P))
}
{ } {
}
{ } {
}
{ } {
}
}
{}{}{.}
} {Pemetaan singgung dan proyeksi vertikal keduanya linear. Menurut (1), ketergantungan pada $F$ juga linear dan selaras dengan perkalian oleh fungsi.
}

}



\inputfaktbeweis
{Differenzierbare Hyperfläche/Zweites Tangentialbündel/Induzierter Zusammenhang/Vertikale Ableitung/Fakt}
{Lema}
{}
{

\faktsituation {Misalkan

\mavergleichskette
{\vergleichskette
{ Y
}
{ \subseteq }{ W
}
{ \subseteq }{ \R^n
}
{    }{
}
{   }{
}
}
{}{}{,}
dengan $W$
\definitionsverweis {terbuka}{}{,}
dan himpunan pertama adalah
\definitionsverweis {hipermuka terdiferensialkan}{}{.}
Misalkan pula
\maabbdisp {F} { Y} { \R^n
} {}
adalah medan vektor tangensial.}
\faktfolgerung {Maka
\definitionsverweis {turunan vertikal}{}{}
$\nabla_F$ yang ditentukan oleh
\definitionsverweis {koneksi}{}{}
pada
Catatan 24.8
berimpit dengan
\definitionsverweis {turunan kovarian}{}{.}}
\faktzusatz {}
\faktzusatz {}

}
{

Untuk penampang diferensiabel
\maabb {s} { Y } { TY
} {}
turunan vertikal koneksi tersebut searah $F$ diberikan oleh komposisi

\mathdisp {Y \stackrel{F}{\longrightarrow} TY \stackrel{T(s)}{\longrightarrow} TTY  \stackrel{  \pi_{\text{vert} } }{\longrightarrow} TY} { . }

Pemetaan $T(s)$ mengirim titik
\mathl{(P,u)}{} ke
\mathl{(P,u, s(P), \left(  D s   \right)_{ P } \left(  u   \right))}{} dalam notasi
Catatan 24.8,
kemudian proyeksi vertikal mengirimnya ke
\mathl{(P, \pi_P  \left(  \left(  D s   \right)_{ P } \left(  u   \right)  \right) )}{.}
Inilah tepatnya rumus dalam definisi turunan kovarian.

}
